\chapter{Persamaan Diferensial}\label{ch:g12:diffeq}

\omterm{def:g12:diffeq:ode}{Persamaan diferensial} mengaitkan sebuah \omterm{def:g11:func:function}{fungsi} dengan \omterm{def:g12:deriv:derivative}{turunannya}. Fisika,
kimia, biologi dan ekonomi menyatakan hukumnya dalam bentuk ini: laju
pendinginan sebuah benda, peluruhan inti radioaktif, pertumbuhan sebuah
populasi, semuanya pernyataan tentang $y'$. Bab ini menyelesaikan secara
tuntas \omterm{def:g10:algebra:equation}{persamaan} linear tingkat satu dengan koefisien tetap.

\section{Persamaan \texorpdfstring{$y' = ay$}{y' = ay}}

\begin{definition}[Persamaan diferensial]\label{def:g12:diffeq:ode}
Suatu \emph{persamaan diferensial}\index{persamaan diferensial} adalah \omterm{def:g10:algebra:equation}{persamaan}
yang bilangan tak diketahuinya berupa \omterm{def:g11:func:function}{fungsi} $y$, dan yang memuat $y$ beserta
\omterm{def:g12:deriv:derivative}{turunannya}. Suatu \emph{penyelesaian} pada selang $I$ adalah \omterm{def:g11:func:function}{fungsi} yang \omterm{def:g12:deriv:derivative}{dapat
diturunkan} dan memenuhi \omterm{def:g10:algebra:equation}{persamaan} itu di setiap titik $I$.
\end{definition}

\begin{theorem}[Penyelesaian $y' = ay$]\label{thm:g12:diffeq:homogeneous}
Misalkan $a \in \R$. Penyelesaian pada $\R$ dari \omterm{def:g10:algebra:equation}{persamaan} $y' = ay$ tepat
berupa \omterm{def:g11:func:function}{fungsi}
\[
y(x) = C\,\eu^{ax}, \qquad C \in \R .
\]
Untuk setiap pasangan $(x_0, y_0)$ ada tepat satu penyelesaian dengan
$y(x_0) = y_0$.
\end{theorem}

\begin{proof}
Setiap $y = C\eu^{ax}$ memenuhi $y' = Ca\,\eu^{ax} = ay$. Sebaliknya, misalkan
$y$ sebarang penyelesaian dan ambil $z(x) = y(x)\,\eu^{-ax}$. Maka
\[
z'(x) = y'(x)\,\eu^{-ax} - a\,y(x)\,\eu^{-ax}
= \bigl(y'(x) - a y(x)\bigr)\eu^{-ax} = 0,
\]
sehingga $z$ bernilai tetap, sebutlah $C$, dan $y(x) = C\eu^{ax}$. Syarat awal
$y(x_0) = y_0$ memaksa $C = y_0 \eu^{-a x_0}$ secara tunggal.
\end{proof}

\begin{remark}
\omterm{thm:g12:exp:existence}{Fungsi eksponen} dengan demikian \emph{dicirikan} oleh \omterm{def:g12:diffeq:ode}{persamaan diferensial}
yang paling sederhana dari semuanya: pertumbuhan yang sebanding dengan
besarnya. Itulah sebabnya ia muncul di mana-mana di alam.
\end{remark}

\begin{omfigure}
\begin{tikzpicture}[scale=1.05]
  \foreach \x in {-2.8,-2.4,...,2.9}
    \foreach \y in {-1.8,-1.4,...,1.9}
      \draw[gray!60, thin] (\x,\y) ++({atan(\y)}:0.11)
        -- ++({atan(\y)+180}:0.22);
  \draw[-Stealth] (-3.2,0) -- (3.4,0) node[below, font=\small] {$x$};
  \draw[-Stealth] (0,-2.2) -- (0,2.4) node[left, font=\small] {$y$};
  \draw[omDef, thick] plot[domain=-3:0.69, samples=60] (\x, {exp(\x)});
  \draw[omThm, thick] plot[domain=-3:1.89, samples=60] (\x, {0.3*exp(\x)});
  \draw[omProp, thick] plot[domain=-3:1.38, samples=60] (\x, {-0.5*exp(\x)});
\end{tikzpicture}

{\small \omterm{def:g10:algebra:equation}{Persamaan} $y' = y$ menetapkan sebuah \omterm{def:g10:reffunc:affine}{gradien} di setiap titik
bidangnya (ruas kelabu). Penyelesaiannya $C\eu^{x}$ --- di sini $C = 1$ (biru),
$C = 0.3$ (merah), $C = -0.5$ (jingga) --- tepat berupa kurva yang
mengikuti medan \omterm{def:g10:reffunc:affine}{gradien} itu.}
\end{omfigure}

\begin{example}[Peluruhan radioaktif]\label{ex:g12:diffeq:decay}
Suatu banyaknya bahan radioaktif memenuhi $N' = -\lambda N$ dengan $\lambda > 0$,
sehingga $N(t) = N_0\,\eu^{-\lambda t}$; lihat \cref{ex:g12:exp:halflife} untuk
waktu paruhnya.
\end{example}

\section{Persamaan \texorpdfstring{$y' = ay + b$}{y' = ay + b}}

\begin{theorem}[Penyelesaian $y' = ay + b$]\label{thm:g12:diffeq:affine}
Misalkan $a \neq 0$ dan $b \in \R$. Penyelesaian pada $\R$ dari $y' = ay + b$
tepat berupa \omterm{def:g11:func:function}{fungsi}
\[
y(x) = C\,\eu^{ax} - \frac{b}{a}, \qquad C \in \R .
\]
\omterm{def:g11:func:function}{Fungsi} tetap $y_p = -\frac ba$ disebut \emph{penyelesaian setimbang}.
Untuk setiap $(x_0, y_0)$ ada tepat satu penyelesaian dengan $y(x_0) = y_0$; jika
$a < 0$, setiap penyelesaian menuju kesetimbangan $-\frac ba$ ketika
$x \to +\infty$.
\end{theorem}

\begin{proof}
\omterm{def:g11:func:function}{Fungsi} tetap $y_p = -\frac{b}{a}$ memenuhi $y_p' = 0 = a y_p + b$. Sekarang $y$
adalah penyelesaian jika dan hanya jika
\[
(y - y_p)' = y' = ay + b = a(y - y_p) + \underbrace{a y_p + b}_{=\,0}
= a (y - y_p),
\]
yaitu jika dan hanya jika $z = y - y_p$ memenuhi $z' = az$. Menurut
\cref{thm:g12:diffeq:homogeneous}, $z = C\eu^{ax}$, dan dari situ rumusnya,
keberadaannya serta ketunggalannya. Jika $a < 0$, maka $\eu^{ax} \to 0$ ketika
$x \to +\infty$, sehingga $y(x) \to -\frac ba$.
\end{proof}

\begin{method}[Menyelesaikan $y' = ay + b$ dengan syarat awal]\label{met:g12:diffeq:solve}
\begin{enumerate}
  \item Carilah penyelesaian setimbangnya $y_p = -\frac{b}{a}$ (selesaikan $y' = 0$).
  \item Tulislah penyelesaian umumnya $y = C\eu^{ax} + y_p$.
  \item Tentukan $C$ dari syarat awalnya.
  \item Periksalah perilaku jangka panjangnya dengan naluri fisikanya (apakah
        penyelesaiannya menuju kesetimbangan?).
\end{enumerate}
Siasat yang sama --- \emph{satu penyelesaian khusus ditambah penyelesaian umum
\omterm{def:g10:algebra:equation}{persamaan} homogennya} --- meluas ke $y' = ay + f(x)$: lihat
\cref{prop:g12:diffeq:general}.
\end{method}

\begin{example}[Hukum pendinginan Newton]\label{ex:g12:diffeq:cooling}
Secangkir kopi bersuhu $80\,^\circ$C berada di ruangan bersuhu $20\,^\circ$C. Hukum
Newton menyatakan bahwa suhunya $T$ memenuhi $T' = -k(T - 20)$ untuk suatu $k > 0$,
\emph{yaitu} $T' = -kT + 20k$. Kesetimbangannya $20$, dan
$T(t) = 20 + 60\,\eu^{-kt}$: kopinya mendingin secepat eksponensial menuju suhu
ruangannya.
\end{example}

\begin{omfigure}
\begin{tikzpicture}
\begin{axis}[omaxis,
  width=9.2cm, height=5.8cm,
  xlabel={$t$}, ylabel={$T$},
  xmin=0, xmax=10.5, ymin=0, ymax=88,
  xtick={5,10}, ytick={20,50,80}]
  \addplot[omProp, dashed, domain=0:10.4] {20};
  \addplot[omDef, thick, domain=0:10.2, samples=60] {20 + 60*exp(-0.5*x)};
  \addplot[omDef!50, thick, domain=0:10.2, samples=60] {20 + 30*exp(-0.5*x)};
  \addplot[omDef!50, thick, domain=0:10.2, samples=60] {20 - 15*exp(-0.5*x)};
  \node[omProp, font=\small, anchor=south west] at (axis cs:6.6,20.8)
    {setimbang};
\end{axis}
\end{tikzpicture}

{\small Berapa pun suhu awalnya, semua penyelesaian
$T' = -k(T - 20)$ menuju kesetimbangan $T = 20$ secara eksponensial.}
\end{omfigure}

\section{Persamaan \texorpdfstring{$y' = ay + f(x)$}{y' = ay + f(x)}}

\begin{proposition}[Susunan himpunan penyelesaiannya]\label{prop:g12:diffeq:general}
Misalkan $f$ \omterm{def:g12:limcont:continuity}{kontinu} pada selang $I$ dan $y_p$ satu penyelesaian khusus
\omterm{def:g10:algebra:equation}{persamaan}
\[
y' = ay + f(x)
\]
pada $I$. Maka penyelesaiannya pada $I$ tepat berupa \omterm{def:g11:func:function}{fungsi}
$y = C\eu^{ax} + y_p$, $C \in \R$.
\end{proposition}

\begin{proof}
Seperti pada \cref{thm:g12:diffeq:affine}: $y$ adalah penyelesaian jika dan hanya
jika $z = y - y_p$ memenuhi
$z' = (ay + f) - (ay_p + f) = az$, jika dan hanya jika $z = C\eu^{ax}$.
\end{proof}

\begin{method}[Menerka satu penyelesaian khusus]\label{met:g12:diffeq:particular}
Carilah satu penyelesaian khusus \emph{yang sebentuk} dengan $f$:
\begin{itemize}
  \item $f$ suku banyak berderajat $n$: cobalah suku banyak berderajat $n$;
  \item $f(x) = \alpha\,\eu^{kx}$ dengan $k \neq a$: cobalah $y_p = \beta\,\eu^{kx}$;
  \item $f(x) = \alpha\,\eu^{ax}$ (kasus beresonansi): cobalah
        $y_p = \beta x\,\eu^{ax}$.
\end{itemize}
Substitusikan ke dalam \omterm{def:g10:algebra:equation}{persamaannya} lalu samakan koefisiennya.
\end{method}

\begin{example}\label{ex:g12:diffeq:particular}
Selesaikan $y' = 2y + 4x$. Cobalah $y_p = \alpha x + \beta$:
$\alpha = 2(\alpha x + \beta) + 4x$ untuk semua $x$ memaksa
$2\alpha + 4 = 0$ dan $\alpha = 2\beta$, jadi $\alpha = -2$, $\beta = -1$.
Penyelesaian umumnya: $y = C\eu^{2x} - 2x - 1$.
\end{example}

\section{Latihan}

\begin{exercise}[$\star$]\label{exo:g12:diffeq:1}
Selesaikan pada $\R$:
(a) $y' = 3y$ dengan $y(0) = 2$;\quad
(b) $2y' + y = 0$ dengan $y(0) = -1$;\quad
(c) $y' = -y + 5$ dengan $y(0) = 0$.
\end{exercise}

\begin{exercise}[$\star$]\label{exo:g12:diffeq:2}
Sebuah populasi bakteri tumbuh dengan laju yang sebanding dengan besarnya, dan
berlipat dua setiap $3$ jam. Tulislah \omterm{def:g12:diffeq:ode}{persamaan diferensial} yang dipenuhi
populasinya $N(t)$ lalu tentukan tetapan pembandingnya.
\end{exercise}

\begin{exercise}[$\star$]\label{exo:g12:diffeq:3}
Periksalah bahwa $y_p(x) = x\,\eu^{x}$ adalah penyelesaian $y' = y + \eu^x$, lalu
berikan semua penyelesaiannya pada $\R$.
\end{exercise}

\begin{exercise}[$\star\star$]\label{exo:g12:diffeq:4}
Selesaikan $y' = -2y + \eu^{x}$ dengan $y(0) = 1$.
(Petunjuk: carilah satu penyelesaian khusus berbentuk $\beta\,\eu^{x}$.)
\end{exercise}

\begin{exercise}[$\star\star$]\label{exo:g12:diffeq:5}
Karbon-14 meluruh dengan waktu paruh $5730$ tahun. Sebuah contoh arkeologis
memuat $60\%$ karbon-14 milik makhluk hidup. Taksirlah umurnya.
\end{exercise}

\begin{exercise}[$\star\star$]\label{exo:g12:diffeq:6}
Sebuah tangki memuat $100$ L air murni. Air garam yang memuat $0.2$ kg garam
tiap liternya mengalir masuk dengan laju $5$ L/menit, dan campurannya (yang
dijaga seragam) mengalir keluar dengan laju yang sama. Misalkan $m(t)$ massa
garam di dalam tangki pada waktu $t$ (dalam menit).
\begin{enumerate}
  \item Berilah alasan bahwa $m' = 1 - \dfrac{m}{20}$.
  \item Selesaikan, lalu tentukan limit $m(t)$ ketika $t \to +\infty$.
        Tafsirkan hasilnya.
\end{enumerate}
\end{exercise}

\begin{exercise}[$\star\star$]\label{exo:g12:diffeq:7}
Seorang penerjun payung bermassa $80$ kg jatuh di bawah pengaruh gravitasi
($g = 9.8\ \text{m/s}^2$) dan hambatan udara yang sebanding dengan lajunya,
sehingga lajunya memenuhi $v' = g - \frac{k}{m}v$ dengan $k = 16$ kg/s.
\begin{enumerate}
  \item Selesaikan \omterm{def:g10:algebra:equation}{persamaannya} dengan $v(0) = 0$.
  \item Hitunglah \emph{kecepatan terminalnya} $\lim_{t\to+\infty} v(t)$, dan
        waktu yang diperlukan untuk mencapai $95\%$ darinya.
\end{enumerate}
\end{exercise}

\begin{exercise}[$\star\star\star$]\label{exo:g12:diffeq:8}
\emph{(\omterm{def:g10:algebra:equation}{Persamaan} logistik.)} Sebuah populasi $y(t) \in \intoo{0}{1}$ (sebagai
pecahan dari populasi maksimalnya) memenuhi
\[
y' = y(1 - y).
\]
\begin{enumerate}
  \item Misalkan $z = \dfrac{1}{y}$. Tunjukkan bahwa $z$ memenuhi \omterm{def:g10:algebra:equation}{persamaan}
        linear $z' = -z + 1$.
  \item Selesaikan untuk $z$, lalu untuk $y$, dengan $y(0) = \frac{1}{10}$.
  \item Tunjukkan bahwa $y(t) \to 1$ ketika $t \to +\infty$ lalu lukiskan bentuk
        kurva penyelesaiannya.
\end{enumerate}
\end{exercise}

\begin{exercise}[$\star\star\star$]\label{exo:g12:diffeq:9}
Misalkan $y$ suatu penyelesaian $y' = ay + b$ dan $u_n = y(n)$ untuk $n \in \N$.
Tunjukkan bahwa $(u_n)$ memenuhi rekurensi aritmetika-geometri
$u_{n+1} = q u_n + r$, lalu nyatakan $q$ dan $r$ dengan $a$ dan $b$. Syarat
$\abs{q} < 1$ bersesuaian dengan apa bagi \omterm{def:g12:diffeq:ode}{persamaan
diferensialnya}?
\end{exercise}

\section{Soal: Jam di dalam benda}

\begin{problem}[{Soal akhir pekan --- karbon-14 menanggali gua,
hukum pendinginan Newton menakar waktu kejahatan, dan satu persamaan kecil
mengenakan empat kostum}]
\label{pb:g12:diffeq:1}
\omterm{def:g12:diffeq:ode}{Persamaan diferensial} adalah hukum perubahan; menyelesaikannya mengubah
hukum itu menjadi sebuah \emph{jam}. \omterm{def:g10:algebra:equation}{Persamaan} mungil yang sama,
$y' = ay + b$ (\cref{thm:g12:diffeq:affine}), berdetak di dalam
arang prasejarah, kopi yang mendingin, tetes hujan yang jatuh dan
infus rumah sakit --- dan membaca jam itulah urusan
soal ini: ia menanggali lukisan Lascaux, memastikan saat kematian,
lalu mengaudit anggapannya sendiri seperti ilmuwan yang baik.

\medskip
\noindent\textbf{Bagian I --- Kelancaran.}
\begin{enumerate}
  \item Selesaikan $y' = 3y$, $y(0) = 2$; lalu
        $y' = -2y + 6$, $y(0) = 0$
        (\cref{met:g12:diffeq:solve}).
  \item Turunkan kembali muslihat ketunggalan di balik
        \cref{thm:g12:diffeq:homogeneous}: jika $y' = ay$, hitunglah
        \omterm{def:g12:deriv:derivative}{turunan} $y(t)\,\eu^{-at}$ lalu simpulkan bahwa
        \emph{setiap} penyelesaiannya berbentuk $C\eu^{at}$.
  \item Untuk $y' = -2y + 6$: carilah penyelesaian setimbangnya, lalu
        lukiskan nasib setiap penyelesaian lainnya ketika
        $t \to +\infty$. (\omterm{pb:g10:functions:1}{Titik tetap} rekurensi pada
        \cref{pb:g11:seq:1}, yang kini menjadi sinambung.)
  \item Untuk hukum peluruhan $y' = -\frac{y}{\tau}$: tunjukkan bahwa
        waktu paruhnya $t_{1/2} = \tau \ln 2$.
  \item Lukiskan (atau uraikan) keluarga penyelesaian
        $y' = -2y + 6$: apa yang dilakukan penyelesaian yang bermula di atas
        $3$? Yang di bawah $3$? Yang tepat di $3$?
\end{enumerate}

\medskip
\noindent\textbf{Bagian II --- Karbon-14.}
Jaringan hidup mempertahankan nisbah karbon-14 radioaktif yang
tetap; pada saat mati asupannya berhenti dan persediaannya meluruh:
$N' = -\lambda N$, dengan waktu paruh $5\,730$ tahun.
\begin{enumerate}[resume]
  \item Hitunglah $\lambda$ (tiap tahunnya).
  \item Sebuah tulang menyisakan $20\,\%$ nisbah makhluk hidupnya: berapa
        umurnya?
  \item Arang dari gua berlukis di Lascaux menyisakan
        kira-kira $15\,\%$: tanggalilah lukisannya.
  \item Manusia es Ötzi, yang ditemukan di sebuah gletser Alpen, terukur
        kira-kira $53\,\%$: tanggalilah dia (para arkeolog menyebut kira-kira
        $5\,300$ tahun --- bagaimana hasilmu?).
  \item Mengapa karbon-14 tidak dapat menanggali dinosaurus? Hitunglah
        pecahan yang tersisa sesudah sepuluh waktu paruh, nyatakan
        pecahannya sesudah $65$ juta tahun sebagai pangkat $2$,
        lalu simpulkan.
  \item Batang galatnya: jika $53\,\%$ milik Ötzi hanya diketahui sampai
        $\pm 1\,\%$, hitunglah rentang umurnya. Ketelitian laboratorium
        seperti apa membeli ketelitian penanggalan seperti apa?
\end{enumerate}

\medskip
\noindent\textbf{Bagian III --- Pendinginan Newton, dan sebuah kejahatan.}
Sebuah benda bersuhu $T$ di lingkungan bersuhu tetap $T_a$
mendingin menurut $T' = -k\,(T - T_a)$.
\begin{enumerate}[resume]
  \item Selesaikan \omterm{def:g10:algebra:equation}{persamaannya} (substitusikan $z = T - T_a$):
        \[
        T(t) = T_a + (T_0 - T_a)\,\eu^{-kt} .
        \]
  \item Kopi yang dituang pada $90\,^\circ$C di ruangan $20\,^\circ$C
        terbaca $70\,^\circ$C sesudah $5$ menit. Tentukan $k$,
        lalu lama menunggu sampai suhunya enak diminum,
        $55\,^\circ$C.
  \item Pertanyaan susunya: agar dapat meminum kopi sehangat mungkin
        dalam sepuluh menit, haruskah susu dinginnya dimasukkan sekarang atau
        pada saat terakhir? Jawablah dengan \omterm{def:g10:algebra:equation}{persamaannya} (apa yang
        dilakukan penambahan susu terhadap jurang $T - T_a$, dan bagaimana
        jurang itu menggerakkan kehilangan panasnya?).
  \item Forensik: sesosok jenazah ditemukan tengah malam pada
        $30\,^\circ$C di ruangan $20\,^\circ$C; satu jam kemudian
        terbaca $28\,^\circ$C. Dengan menganggap $37\,^\circ$C pada
        saat kematiannya, perolehlah $k$ dari kedua pengukuran itu,
        lalu hitunglah saat kematiannya.
  \item Auditlah jam sang dokter forensik: sebutkan tiga pelanggaran nyata
        terhadap anggapan modelnya dan ke arah mana masing-masing
        akan membiaskan taksiran saat kematiannya.
\end{enumerate}

\medskip
\noindent\textbf{Bagian IV --- Empat kostum dan sebuah kurva S.}
\begin{enumerate}[resume]
  \item \omterm{def:g10:algebra:equation}{Persamaan} logistik pada \cref{exo:g12:diffeq:8},
        $y' = y(1 - y)$ dengan $y(0) = 0.1$: dengan memakai substitusi
        latihan itu, turunkanlah
        $y(t) = \dfrac{1}{1 + 9\eu^{-t}}$, lalu carilah waktu
        \omterm{def:g12:deriv:inflection}{titik beloknya} $y = \frac12$ --- tanggal milik
        ahli epidemiologi pada \cref{pb:g12:deriv:1}, yang kini
        dapat dihitung.
  \item Setetes hujan menuruti $v' = 10 - \frac v2$ (gravitasi dikurangi
        hambatan), $v(0) = 0$. Carilah kecepatan terminalnya, selesaikan
        untuk $v(t)$, lalu hitunglah kapan tetesnya mencapai
        $95\,\%$ laju terminalnya.
  \item Sebuah infus rumah sakit menyalurkan obat dengan laju tetap sementara
        tubuh membuangnya secara sebanding:
        $c' = 4 - 0.5c$, $c(0) = 0$. Carilah kepekatan pada keadaan
        mapannya dan waktu untuk mencapai setengahnya. Sebutkan
        dalam satu kalimat kesejajarannya dengan rekurensi pinjaman pada
        \cref{exo:g12:diffeq:9}.
  \item Penutup --- satu \omterm{def:g10:algebra:equation}{persamaan}, empat kostum: peluruhan, pendinginan,
        jatuh, penakaran obat, semuanya $y' = ay + b$ dengan tanda dan nama
        yang berbeda. Lafalkanlah gelung permodelan yang dijalankan soal ini
        empat kali (hukum $\to$ \omterm{def:g10:algebra:equation}{persamaan} $\to$ penyelesaian $\to$
        penakaran $\to$ ramalan $\to$ audit), lalu sebutkan
        gejala yang menuntut $y''$ --- dan soal mana yang
        sudah menjumpainya.

\end{enumerate}
\end{problem}
